% VI31 detailed challenge solutions -- VI/24+ proof-depth standard.

% VI31-DETAILED-CHALLENGE-01
\subsection*{Challenge VI.31.C1 solution}
\begin{solution}
Work in
\[
A=k[x_1,x_2,x_3,x_4].
\]

For codimension \(1\), use the three-equation presentation
\[
I_1=(x_1,x_1x_2,x_1x_3).
\]
Since the last two generators already lie in \((x_1)\),
\[
I_1=(x_1),
\qquad
\sqrt{I_1}=(x_1).
\]
Thus
\[
V(I_1)=V(x_1)
\]
has codimension \(1\).  A realizing prime chain is
\[
(0)\subsetneq(x_1).
\]

For codimension \(2\), take
\[
I_2=(x_1,x_2,x_1x_3).
\]
Again the third equation is redundant:
\[
I_2=(x_1,x_2),
\qquad
\sqrt{I_2}=(x_1,x_2).
\]
Hence
\[
V(I_2)=V(x_1,x_2)
\]
has codimension \(2\), realized by
\[
(0)\subsetneq(x_1)\subsetneq(x_1,x_2).
\]

For codimension \(3\), take
\[
I_3=(x_1,x_2,x_3).
\]
This ideal is prime, because
\[
A/I_3\cong k[x_4]
\]
is a domain.  The chain
\[
(0)\subsetneq(x_1)\subsetneq(x_1,x_2)\subsetneq(x_1,x_2,x_3)
\]
has length \(3\), so
\[
\operatorname{codim}V(I_3)=3.
\]

Thus three equations can describe loci of codimension \(1\), \(2\), or \(3\):
\[
\boxed{
\operatorname{codim}V(I_1)=1,\quad
\operatorname{codim}V(I_2)=2,\quad
\operatorname{codim}V(I_3)=3.
}
\]
The radical ideal and its prime chains, not the displayed equation count, control codimension.
\end{solution}

% VI31-DETAILED-CHALLENGE-02
\subsection*{Challenge VI.31.C2 solution}
\begin{solution}
Let \(Y\subseteq X\) be irreducible closed.  The scheme-theoretic generic-point property supplies a unique point
\[
\eta_Y\in Y
\]
such that
\[
\overline{\{\eta_Y\}}=Y.
\]
This is the first place where sobriety enters: without a generic point, there would be no single point whose
specialization order represents the whole irreducible closed subset.

A chain contributing to codimension is
\[
Y=Z_n\subsetneq Z_{n-1}\subsetneq\cdots\subsetneq Z_0,
\]
with each \(Z_i\) irreducible and closed.  Let \(\eta_i\) be the generic point of \(Z_i\).  Since
\[
Z_{i+1}\subsetneq Z_i,
\]
we have
\[
\eta_i\in\overline{\{\eta_{i+1}\}},
\]
so \(\eta_{i+1}\) is a specialization of \(\eta_i\), or equivalently \(\eta_i\) is a strict generalization of
\(\eta_{i+1}\).  Thus the closed-set chain becomes a strict generalization chain ending at \(\eta_Y\).

Choose an affine neighborhood
\[
U=\Spec A
\]
of \(\eta_Y\).  Open subsets are stable under generalization, so every point \(\eta_i\) in the chain also lies in
\(U\).  If \(\eta_Y\) corresponds to \(\mathfrak p\), the generalizations correspond exactly to primes
\[
\mathfrak p_0\subsetneq\mathfrak p_1\subsetneq\cdots\subsetneq\mathfrak p_n=\mathfrak p.
\]
Therefore codimension is the supremum of lengths of prime chains ending at \(\mathfrak p\):
\[
\operatorname{codim}(Y,X)=\operatorname{ht}(\mathfrak p).
\]

Finally
\[
\mathcal O_{X,\eta_Y}\cong A_{\mathfrak p},
\]
and prime ideals of \(A_{\mathfrak p}\) correspond to primes of \(A\) contained in \(\mathfrak p\).  Hence
\[
\dim A_{\mathfrak p}=\operatorname{ht}(\mathfrak p).
\]
Combining the identifications,
\[
\boxed{
\operatorname{codim}(Y,X)=\dim\mathcal O_{X,\eta_Y}.
}
\]

The sobriety/generic-point property is used twice conceptually: first to replace \(Y\) and each \(Z_i\) by unique
points, and then to translate inclusion of irreducible closed subsets into the specialization order on those
points.
\end{solution}

% VI31-DETAILED-CHALLENGE-03
\subsection*{Challenge VI.31.C3 solution}
\begin{solution}
There are three logically separate ingredients in the variety proof.

First, the chain-theoretic identity
\[
\operatorname{codim}(Y,X)
=
\dim\mathcal O_{X,\eta_Y}
\]
is very general.  It uses only the scheme specialization order and the affine height dictionary.  It does not
need finite type over a field.

Second, the affine formula
\[
\dim A
=
\operatorname{ht}(\mathfrak p)+\dim(A/\mathfrak p)
\]
is not a formal property of every ring.  In the finitely generated \(k\)-domain setting it follows from the
finite-type dimension theory and is exactly what converts relative height into a difference of global
dimensions.

Third, VI/30 supplies open-dimension invariance for varieties, so an affine chart meeting \(Y\) has the same
dimension as \(X\), and its intersection with \(Y\) has the same dimension as \(Y\).

For a variety, all three ingredients are available and yield
\[
\boxed{
\dim Y+\operatorname{codim}(Y,X)=\dim X.
}
\]

For the non-equidimensional disjoint union
\[
X=\Spec k\sqcup\mathbb A^1_k,
\qquad
Y=\Spec k,
\]
the chain-theoretic codimension formula still works:
\[
\operatorname{codim}(Y,X)=0.
\]
But \(Y\) lies on a zero-dimensional irreducible component while the global dimension of \(X\) is controlled by
the other one-dimensional component.  Thus
\[
0+0\ne1.
\]
The failure is global component bookkeeping, not a failure of the height dictionary.

A general Noetherian affine scheme can also fail without suitable equidimensional/catenary hypotheses.  A simple
affine example is
\[
A=k[t]\times k[x,y].
\]
Then
\[
X=\Spec A
\]
is Noetherian with components of dimensions \(1\) and \(2\).  If \(Y=\Spec k[t]\) is the first component, then
\[
\dim Y=1,\qquad
\operatorname{codim}(Y,X)=0,\qquad
\dim X=2,
\]
so additivity fails.

For irreducible Noetherian schemes, the height-plus-quotient formula also requires extra hypotheses.
Noncatenary Noetherian domains give the deeper warning.  VI/31 therefore states exact additivity in the safe
finite-type variety setting, not for all Noetherian schemes.
\end{solution}

% VI31-DETAILED-CHALLENGE-04
\subsection*{Challenge VI.31.C4 solution}
\begin{solution}
Assume the products in question are irreducible and have the expected dimensions.  Let
\[
Y_1\subseteq X_1,
\qquad
Y_2\subseteq X_2
\]
be irreducible closed subvarieties.  Then
\[
\dim(X_1\times_kX_2)=\dim X_1+\dim X_2
\]
and
\[
\dim(Y_1\times_kY_2)=\dim Y_1+\dim Y_2.
\]
Applying dimension--codimension additivity,
\[
\begin{aligned}
\operatorname{codim}(Y_1\times Y_2,X_1\times X_2)
&=
\dim(X_1\times X_2)-\dim(Y_1\times Y_2)\\
&=
(\dim X_1-\dim Y_1)+(\dim X_2-\dim Y_2).
\end{aligned}
\]
Therefore
\[
\boxed{
\operatorname{codim}(Y_1\times Y_2,X_1\times X_2)
=
\operatorname{codim}(Y_1,X_1)
+
\operatorname{codim}(Y_2,X_2).
}
\]

For the diagonal, one additional point is essential in the chapter's convention for ``variety'': to discuss
\(\Delta_X\) as a closed subvariety, assume \(X\) is separated over \(k\).  Then
\[
\Delta_X\hookrightarrow X\times_kX
\]
is closed and
\[
\Delta_X\cong X.
\]
If \(d=\dim X\) and the product has expected dimension \(2d\), then
\[
\boxed{
\operatorname{codim}(\Delta_X,X\times_kX)=2d-d=d.
}
\]
For affine varieties the separatedness condition is automatic.
\end{solution}

% VI31-DETAILED-CHALLENGE-05
\subsection*{Challenge VI.31.C5 solution}
\begin{solution}
Let \(X\) be normal, Noetherian, and integral, and let \(x\in X\) be a codimension-one point.  The closure
\[
\overline{\{x\}}
\]
is an irreducible closed subset of codimension one.  By the generic-point local-ring theorem,
\[
\dim\mathcal O_{X,x}=1.
\]
Normality is local, so
\[
\mathcal O_{X,x}
\]
is an integrally closed local domain; Noetherianity is also local.  Thus the codimension-one local ring has the
three structural properties
\[
\boxed{
\text{Noetherian}+\text{local}+\text{integrally closed domain of dimension }1.
}
\]

This is precisely the algebraic configuration from which the later valuation description arises.  VI/31 stops
here: codimension identifies the points and proves their local rings are one-dimensional normal local domains.
The later divisor chapter may then invoke the additional commutative-algebra theorem that turns this
one-dimensional normal local structure into the discrete valuation framework used to measure zeros and poles.

The following statements are therefore deliberately \emph{not} proved in VI/31:
\begin{itemize}
\item that a codimension-one local ring is used as a discrete valuation ring;
\item construction of an order function \(\operatorname{ord}_x:K(X)^\times\to\mathbb Z\);
\item finiteness of the nonzero orders of a rational function;
\item the principal divisor
\[
\operatorname{div}(f)=\sum_{x\in X^{(1)}}\operatorname{ord}_x(f)\,[\overline{\{x\}}];
\]
\item the Weil divisor group and linear equivalence;
\item Cartier divisors and comparison with Weil divisors;
\item divisor class groups.
\end{itemize}
Those belong to the later divisor arc.  The output of VI/31 is only the numerical and local-ring bridge
\[
\text{codimension one}
\Longleftrightarrow
\text{height one}
\Longleftrightarrow
\dim\mathcal O_{X,x}=1.
\]
\end{solution}
